\section{Edge indicators}\label{edg:section}
This section proves that every centered edge indicator is an eigenfunction of
$H$ with eigenvalue $n-1$. It uses only the balance \eqref{eq:balance} of a
regular fiber. Section~\ref{root:section} uses the result for the upper bound
in Theorem~\ref{thm:main}.

\begin{lemma}[Balanced trades]\label{edg:marginal}
Let $a=\{x,w\}$ and $y\notin a$. Then
\[
 E_a\mathbf1[xy]=\tfrac12\bigl(\mathbf1[xy]+\mathbf1[wy]\bigr).
\]
If a pair $e$ equals $a$ or is disjoint from $a$, then
$E_a\mathbf1[e]=\mathbf1[e]$.
\end{lemma}

\begin{proof}
Fix an $a$-fiber. If $y$ is common or adjacent to neither endpoint, then
$\mathbf1[xy]=\mathbf1[wy]$ is constant on the fiber, and both sides agree. If
$y$ is a singleton, then $\mathbf1[xy]+\mathbf1[wy]=1$ on the fiber. By
\eqref{eq:balance}, the set of singletons adjacent to $x$ is a uniform
$q$-subset of the $2q$ singletons, which contains $y$ with probability $1/2$.
So $E_a\mathbf1[xy]=1/2$ on the fiber. The last assertion holds because a
trade at $a$ changes only pairs joining $a$ to a singleton.
\end{proof}

\begin{proposition}[Edge eigenfunctions]\label{edg:eigen}
For every pair $\{x,y\}$,
\[
 H\mathbf1[xy]=(n-1)\mathbf1[xy]-d .
\]
With $p=d/(n-1)$, the function $g_{xy}=\mathbf1[xy]-p$ is nonzero, has mean
zero, and satisfies $Hg_{xy}=(n-1)g_{xy}$.
\end{proposition}

\begin{proof}
By Lemma~\ref{edg:marginal}, $h_a\mathbf1[xy]=0$ unless $a$ meets $\{x,y\}$ in
exactly one vertex. For $w\notin\{x,y\}$ the pairs $\{x,w\}$ and $\{y,w\}$
contribute
\[
 \tfrac12\bigl(\mathbf1[xy]-\mathbf1[wy]\bigr)
 +\tfrac12\bigl(\mathbf1[xy]-\mathbf1[xw]\bigr).
\]
The vertex $y$ has $d-\mathbf1[xy]$ neighbors other than $x$, and $x$ has
$d-\mathbf1[xy]$ neighbors other than $y$. Summing over the $n-2$ choices of
$w$ gives
\[
 H\mathbf1[xy]=(n-2)\mathbf1[xy]-\bigl(d-\mathbf1[xy]\bigr)
 =(n-1)\mathbf1[xy]-d .
\]
Since $H\mathbf1=0$, it follows that $Hg_{xy}=(n-1)\mathbf1[xy]-d=(n-1)g_{xy}$.
By \eqref{eq:density}, $g_{xy}$ has mean zero and takes the two values $1-p$
and $-p$, each with positive probability, so it is nonzero.
\end{proof}

\begin{lemma}[Complement]\label{edg:complement}
Complementation $G\mapsto\overline G$ is a bijection from $\Omega_{n,d}$ onto
$\Omega_{n,n-1-d}$ which maps each $a$-fiber onto an $a$-fiber. Hence the
operators $H$ on the two spaces are unitarily equivalent by a map that
commutes with relabeling, and they have the same eigenvalues on rooted-class
functions.
\end{lemma}

\begin{proof}
Let $a=\{x,w\}$. A vertex $y\notin a$ is a singleton for $G$ exactly when it is
a singleton for $\overline G$, and then it is adjacent in $\overline G$ to the
endpoint it is not adjacent to in $G$. The two graphs have complementary
edges on every pair. So two graphs agree off the pairs joining $a$ to the
singletons exactly when their complements do, and the fibers correspond.
Composition with the bijection is therefore a unitary map
$L^2(\Omega_{n,n-1-d})\to L^2(\Omega_{n,d})$ that intertwines the operators
$E_a$, hence $H$. Complementation commutes with relabeling, so the map
preserves rooted-class functions.
\end{proof}
